\chapter{Project Requirements}
\label{sec:requirements}
This chapter provides a comprehensive specification of the functional and non-functional requirements for the complete MyAIStorybook system across all FYP iterations. The requirements are organized to reflect the modular architecture and user interactions supported by the full platform implementation, encompassing foundational story generation capabilities through advanced interactive features.

\section{Use Case Analysis}
The MyAIStorybook platform supports comprehensive operational workflows encompassing story generation, interactive features, and advanced user engagement capabilities. The following use case descriptions capture the primary functionalities and system behaviors across all FYP iterations.

% SPACE FOR USE CASE DIAGRAM - TO BE ADDED
\vspace{3cm}
\begin{center}
\textit{[Use Case Diagram will be inserted here showing all 19 use cases across text input, voice input, chat interaction, media generation, export features, and administrative functions]}
\end{center}
\vspace{2cm}

\subsection{UC1: Generate Story from Prompt}
\textbf{Primary Actor:} End User (Parent, Educator, Content Creator)\\
\textbf{Preconditions:} System is running; both backend and frontend are accessible\\
\textbf{Postconditions:} Complete illustrated storybook is generated, displayed, and exported as PDF\\
\textbf{Main Success Scenario:}
\begin{enumerate}
    \item User navigates to the landing page and clicks "Start Creating"
    \item User enters a story prompt in the text area input field
    \item User optionally toggles the "Illustrate My Story" checkbox to enable/disable image generation
    \item User clicks the "Weave Magic" button to initiate generation
    \item System validates the prompt through Prompt Agent
    \item System enhances the prompt and classifies it (normal, detailed, invalid, or nonsense)
    \item If invalid or nonsense, system returns error message; otherwise, proceeds
    \item System orchestrates the multi-agent pipeline via Director Agent
    \item Writer Agent generates story structure with 3 scenes using Ollama LLM
    \item Reviewer Agent validates content for age-appropriateness and coherence
    \item Editor Agent polishes the story and generates PDF export
    \item If images enabled, Image Agent generates illustrations for each scene using Stable Diffusion
    \item Evaluation Agent runs in background to collect quality metrics
    \item System saves all outputs to respective directories under generated/
    \item System returns complete story with image URLs to frontend
    \item Frontend displays the story in page-by-page navigation interface
    \item User navigates through pages and views the generated content
\end{enumerate}

\textbf{Alternative Flows:}
\begin{itemize}
    \item 4a. User presses Enter in textarea instead of clicking button → same behavior
    \item 7a. Prompt is invalid → System displays error message, user modifies prompt
    \item 12a. Image generation fails → System continues with text-only story, logs error
\end{itemize}

\subsection{UC2: Generate Story from Voice Input}
\textbf{Primary Actor:} End User (Parent, Educator, Child)\\
\textbf{Preconditions:} System is running; microphone access is granted\\
\textbf{Postconditions:} Complete illustrated storybook is generated from voice input, displayed, and exported as PDF\\
\textbf{Main Success Scenario:}
\begin{enumerate}
    \item User navigates to the landing page and selects voice input mode
    \item User speaks a story prompt into the microphone
    \item System processes voice input through STT Agent using OpenAI Whisper
    \item System transcribes speech to text and validates the prompt
    \item System proceeds with standard story generation pipeline
    \item System generates story with audio narration using TTS Agent
    \item System displays story with synchronized audio playback
    \item User can listen to the story narration while reading
\end{enumerate}

\subsection{UC3: Interactive Story Chat}
\textbf{Primary Actor:} End User (Parent, Educator, Child)\\
\textbf{Preconditions:} System is running; chat mode is enabled\\
\textbf{Postconditions:} Interactive conversation with AI about story content, personalized recommendations provided\\
\textbf{Main Success Scenario:}
\begin{enumerate}
    \item User initiates chat mode from the main interface
    \item System initializes Chatbot Agent with conversation context
    \item User types or speaks a message about story preferences or questions
    \item System processes input through Chatbot Agent using Ollama LLM
    \item System generates contextual response based on conversation history
    \item User continues conversation to refine story requirements
    \item System adapts recommendations based on user preferences
    \item User can generate final story based on conversation insights
\end{enumerate}

\subsection{UC4: Generate Story Images with Character Consistency}
\textbf{Primary Actor:} End User (Content Creator, Educator)\\
\textbf{Preconditions:} Story generation is in progress; image generation is enabled\\
\textbf{Postconditions:} High-quality illustrations with consistent character appearance across scenes\\
\textbf{Main Success Scenario:}
\begin{enumerate}
    \item System generates initial scene illustration using Stable Diffusion
    \item System creates character reference for consistency tracking
    \item For subsequent scenes, system applies character consistency techniques
    \item System uses style transfer to maintain artistic coherence
    \item System applies upscaling and enhancement for high-resolution output
    \item System saves images with character consistency metadata
    \item Frontend displays scenes with consistent character appearance
\end{enumerate}

\subsection{UC5: Generate Story Audio Narration}
\textbf{Primary Actor:} End User (Parent, Educator, Child)\\
\textbf{Preconditions:} Story generation is in progress; audio generation is enabled\\
\textbf{Postconditions:} Natural speech narration with emotion and character voices\\
\textbf{Main Success Scenario:}
\begin{enumerate}
    \item System processes story text through TTS Agent using Coqui TTS
    \item System detects emotional tone and applies appropriate voice style
    \item System generates character-specific voice variations
    \item System adds background music and sound effects
    \item System synchronizes audio with story text and page navigation
    \item System saves audio files with metadata for playback
    \item Frontend provides audio controls for synchronized narration
\end{enumerate}

\subsection{UC6: Export Story to PDF}
\textbf{Primary Actor:} End User (Parent, Educator)\\
\textbf{Preconditions:} A complete story has been generated and displayed\\
\textbf{Postconditions:} PDF document is created and stored in generated/exports directory\\
\textbf{Main Success Scenario:}
\begin{enumerate}
    \item Story generation completes successfully
    \item Editor Agent automatically generates PDF during the editing phase
    \item PDF includes story title, scenes, and embedded images if generated
    \item PDF is saved with timestamped filename in generated/exports/
    \item User can navigate to the exports directory to access the PDF file
\end{enumerate}

\subsection{UC7: Export Story to Audio Book}
\textbf{Primary Actor:} End User (Parent, Educator)\\
\textbf{Preconditions:} A complete story with audio has been generated\\
\textbf{Postconditions:} Audio book file is created with synchronized narration\\
\textbf{Main Success Scenario:}
\begin{enumerate}
    \item System combines all scene audio files into continuous narration
    \item System adds chapter markers and navigation points
    \item System applies audio enhancement and normalization
    \item System exports audio book in standard format (MP3/WAV)
    \item System saves audio book with metadata and cover information
    \item User can download and share the audio book file
\end{enumerate}

\subsection{UC8: Toggle Theme}
\textbf{Primary Actor:} End User (Any)\\
\textbf{Preconditions:} Application is loaded\\
\textbf{Postconditions:} Visual theme is switched and preference maintained during session\\
\textbf{Main Success Scenario:}
\begin{enumerate}
    \item User clicks the theme toggle button in the navigation bar
    \item System switches between light and dark theme using Theme Context
    \item All components re-render with updated theme styles
    \item Theme preference persists throughout the user session
\end{enumerate}

\section{Functional Requirements}
This section details the specific functional capabilities required from each module of the complete MyAIStorybook system across all FYP iterations.

\subsection{Multi-Agent Story Generation Module}
\begin{enumerate}
    \item FR1.1: The system shall accept user prompts ranging from 5 to 500 characters for story generation
    \item FR1.2: The Prompt Agent shall validate user input and classify prompts as normal, detailed, invalid, or nonsense
    \item FR1.3: The Prompt Agent shall enhance simple prompts by adding contextual details while preserving user intent
    \item FR1.4: The Prompt Agent shall detect input language and determine appropriate age group for content
    \item FR1.5: The Director Agent shall orchestrate the sequential execution of Writer, Reviewer, and Editor agents
    \item FR1.6: The Writer Agent shall generate structured stories with title, setting, characters, and scenes using Ollama LLM (llama3.1:8b-instruct-q4\_K\_M)
    \item FR1.7: The Writer Agent shall create exactly 3 scenes per story as configured
    \item FR1.8: The Writer Agent shall retry up to 2 times if JSON parsing fails
    \item FR1.9: The Writer Agent shall develop character backstories and personality traits
    \item FR1.10: The Writer Agent shall generate dialogue and character interactions
    \item FR1.11: The Reviewer Agent shall validate generated content for coherence, age-appropriateness, and narrative consistency
    \item FR1.12: The Reviewer Agent shall assess educational value and learning objectives
    \item FR1.13: The Editor Agent shall refine story text, improve language quality, and ensure proper formatting
    \item FR1.14: The Editor Agent shall generate PDF export automatically during the editing phase
    \item FR1.15: The TTS Agent shall synthesize natural speech from story text using Coqui TTS
    \item FR1.16: The TTS Agent shall apply emotion detection and voice style customization
    \item FR1.17: The STT Agent shall transcribe voice input using OpenAI Whisper
    \item FR1.18: The STT Agent shall support English language voice input processing
    \item FR1.19: The Chatbot Agent shall process conversational input and maintain context
    \item FR1.20: The Chatbot Agent shall provide personalized story recommendations
    \item FR1.21: The Evaluation Agent shall run in background to collect quality metrics without blocking the main response
    \item FR1.22: The system shall output stories in comprehensive JSON format with metadata and media paths
    \item FR1.23: The system shall save intermediate outputs from each agent in dedicated directories
    \item FR1.24: The system shall maintain processing status and emit console logs at each pipeline stage
    \item FR1.25: The system shall handle errors gracefully and return informative status messages
    \item FR1.26: The system shall learn from user preferences and adapt content generation accordingly
\end{enumerate}

\subsection{Image Generation Module}
\begin{enumerate}
    \item FR2.1: The Image Agent shall generate illustrations for each story scene using Stable Diffusion 1.5 with DreamShaper 8 checkpoint
    \item FR2.2: The system shall employ a two-pass generation process: txt2img synthesis followed by img2img high-resolution refinement
    \item FR2.3: The system shall truncate image prompts to 77 tokens or fewer to ensure CLIP model compatibility
    \item FR2.4: The Image Agent shall utilize DPMSolver scheduler with Karras sigmas for optimized sampling
    \item FR2.5: The system shall use GPU acceleration with half-precision (float16) when CUDA is available
    \item FR2.6: The system shall automatically fall back to CPU processing if GPU is unavailable
    \item FR2.7: The system shall save generated images in PNG format with unique, timestamped filenames
    \item FR2.8: The system shall organize images in the generated/images directory
    \item FR2.9: The Image Agent shall apply positive and negative prompts for quality control
    \item FR2.10: The system shall continue story generation even if image generation fails (graceful degradation)
    \item FR2.11: The system shall return image file paths in the story JSON for frontend access
    \item FR2.12: The system shall support optional image generation via user toggle
\end{enumerate}

\subsection{PDF Export and Document Generation Module}
\begin{enumerate}
    \item FR3.1: The Editor Agent shall automatically generate PDF documents from completed stories
    \item FR3.2: The PDF module shall include story title, setting, characters, and all scenes
    \item FR3.3: The PDF module shall embed generated images if available
    \item FR3.4: The system shall save generated PDFs in the generated/exports directory
    \item FR3.5: The PDF filenames shall include sanitized story title and timestamp
    \item FR3.6: The PDF generation shall use ReportLab library for rendering
\end{enumerate}

\subsection{Audio Generation and Speech Processing Module}
\begin{enumerate}
    \item FR5.1: The TTS Agent shall synthesize natural speech from story text using Coqui TTS
    \item FR5.2: The system shall support multiple voice styles and character-specific narration
    \item FR5.3: The system shall detect emotional tone and apply appropriate voice characteristics
    \item FR5.4: The system shall generate background music and sound effects for immersive experience
    \item FR5.5: The system shall synchronize audio playback with story text and page navigation
    \item FR5.6: The STT Agent shall transcribe voice input using OpenAI Whisper with high accuracy
    \item FR5.7: The system shall support multiple languages for both speech synthesis and recognition
    \item FR5.8: The system shall provide voice command processing for hands-free interaction
    \item FR5.9: The system shall export audio books in standard formats (MP3/WAV)
    \item FR5.10: The system shall provide accessibility features for visually impaired users
    \item FR5.11: The system shall optimize audio quality and compression for efficient storage
    \item FR5.12: The system shall support real-time audio processing for interactive features
\end{enumerate}

\subsection{Interactive Chat and Conversational AI Module}
\begin{enumerate}
    \item FR6.1: The Chatbot Agent shall process natural language conversation input
    \item FR6.2: The system shall maintain conversation context and memory across interactions
    \item FR6.3: The system shall provide personalized story recommendations based on user preferences
    \item FR6.4: The system shall support interactive story modification through conversational interface
    \item FR6.5: The system shall generate emotion-aware responses for age-appropriate interactions
    \item FR6.6: The system shall integrate voice chat capabilities combining STT and TTS
    \item FR6.7: The system shall provide educational Q\&A features for learning and engagement
    \item FR6.8: The system shall support story continuation and expansion through chat interaction
    \item FR6.9: The system shall learn from user interactions and adapt conversation styles
    \item FR6.10: The system shall implement safety and content moderation for child-appropriate interactions
    \item FR6.11: The system shall support multi-language conversation capabilities
    \item FR6.12: The system shall provide context-aware dialogue management
\end{enumerate}

\subsection{User Interface and Frontend Module}
\begin{enumerate}
    \item FR7.1: The frontend shall provide a Next.js-based user interface with responsive design
    \item FR7.2: The system shall display a landing page with project introduction and multiple input options
    \item FR7.3: The story input interface shall support text, voice, and chat input modalities
    \item FR7.4: The system shall provide toggles for enabling/disabling image and audio generation
    \item FR7.5: The system shall display custom loading animations during story generation
    \item FR7.6: The story display interface shall present stories with page-by-page navigation
    \item FR7.7: The system shall display story title, scene text, images, and audio controls for each page
    \item FR7.8: The system shall provide real-time chat interface for conversational interaction
    \item FR7.9: The system shall provide theme toggle for switching between light and dark modes
    \item FR7.10: The system shall maintain theme preference throughout the session using Context API
    \item FR7.11: The navigation bar shall provide reset and home navigation functionality
    \item FR7.12: The system shall display user-friendly error messages when generation fails
    \item FR7.13: The frontend shall communicate with the FastAPI backend via POST requests to /api/generate
    \item FR7.14: The system shall serve generated images and audio through static file routes
    \item FR7.15: The system shall provide Progressive Web App capabilities for offline functionality
    \item FR7.16: The system shall support responsive design for accessibility across devices
\end{enumerate}

\section{Non-Functional Requirements}
This section specifies the quality attributes and constraints that define system behavior beyond functional capabilities.

\subsection{Performance}
\begin{enumerate}
    \item NFR1.1: The Writer Agent shall generate story text within 10-15 seconds for typical prompts on systems with adequate CPU
    \item NFR1.2: The Image Agent shall generate illustrations within 30-60 seconds per scene on GPU-enabled systems (NVIDIA GTX 1060 or equivalent)
    \item NFR1.3: The Image Agent shall complete generation within 2-3 minutes per scene on CPU-only systems
    \item NFR1.4: The frontend shall achieve page load times under 2 seconds for the story display interface
    \item NFR1.5: The system shall handle story generation requests sequentially (one at a time in current implementation)
    \item NFR1.6: Static file serving for images shall have response times under 200 milliseconds
\end{enumerate}

\subsection{Reliability}
\begin{enumerate}
    \item NFR2.1: The Writer Agent shall successfully generate valid JSON for 90\% of reasonable prompts
    \item NFR2.2: The Writer Agent shall retry up to 2 times on JSON parsing failures before returning error
    \item NFR2.3: The system shall continue operation and deliver text-only stories when image generation fails
    \item NFR2.4: The system shall preserve generated content in persistent file storage
    \item NFR2.5: The system shall log all errors to console with sufficient detail for troubleshooting
    \item NFR2.6: The system shall validate all story outputs against Pydantic schema before returning to frontend
\end{enumerate}

\subsection{Usability}
\begin{enumerate}
    \item NFR3.1: Users with no technical background shall be able to generate a story within 3 interactions (navigate, enter prompt, click button)
    \item NFR3.2: The user interface shall provide clear visual feedback during all operations
    \item NFR3.3: The system shall display informative error messages that explain what went wrong
    \item NFR3.4: The story input interface shall support Enter key submission for convenience
    \item NFR3.5: The system shall provide visual loading indicators during generation
    \item NFR3.6: Text elements shall maintain readability with appropriate font sizes and contrast ratios in both light and dark themes
\end{enumerate}

\subsection{Security}
\begin{enumerate}
    \item NFR4.1: All AI processing shall occur locally without transmitting user data to external servers
    \item NFR4.2: The system shall sanitize user inputs for safe filename generation
    \item NFR4.3: File system operations shall prevent directory traversal vulnerabilities
    \item NFR4.4: The system shall not store personally identifiable information
    \item NFR4.5: CORS policies shall be configurable for deployment (currently permissive for development)
\end{enumerate}

\subsection{Maintainability}
\begin{enumerate}
    \item NFR5.1: The codebase shall follow modular design principles with clear separation of concerns
    \item NFR5.2: Each agent module shall be independently implementable and testable
    \item NFR5.3: Code shall include inline comments and docstrings for key functions
    \item NFR5.4: The system shall provide console logging for debugging and monitoring
    \item NFR5.5: The frontend shall use TypeScript for type safety and better maintainability
\end{enumerate}

\subsection{Portability}
\begin{enumerate}
    \item NFR6.1: The backend shall operate on Windows 10/11 systems (primary development platform)
    \item NFR6.2: The system shall function correctly on both CPU-only and GPU-enabled hardware
    \item NFR6.3: The frontend shall be compatible with modern web browsers (Chrome, Firefox, Edge, Safari)
    \item NFR6.4: The system shall clearly document hardware requirements and GPU dependencies
\end{enumerate}

\subsection{Scalability Considerations}
\begin{enumerate}
    \item NFR7.1: The architecture shall support addition of new agent types without major refactoring
    \item NFR7.2: The file-based storage structure shall support migration to database systems in future iterations
    \item NFR7.3: The image generation module shall support alternative diffusion models through configuration changes
    \item NFR7.4: The modular architecture shall facilitate addition of new features in subsequent iterations
\end{enumerate}

\section{Domain Model}
The MyAIStorybook domain comprises several core entities that represent the business logic and data structures of the system. The following domain model illustrates the relationships between key entities.

% SPACE FOR DOMAIN MODEL DIAGRAM - TO BE ADDED
\vspace{3cm}
\begin{center}
\textit{[Domain Model / Class Diagram will be inserted here showing Story, Scene, Agent hierarchy, and their relationships]}
\end{center}
\vspace{3cm}

\subsection{Core Domain Entities}

\textbf{Story}\\
Represents a complete generated storybook with metadata, narrative content, and visual elements.\\
\textit{Attributes:} title, setting, characters (list), scenes (list), meta (optional)

\textbf{Scene}\\
Represents an individual page or segment of the story with associated text and optional illustration.\\
\textit{Attributes:} scene\_number, text, image\_description (optional), image\_path (optional)

\textbf{Agent (Abstract Base)}\\
Abstract representation of specialized processing components in the multi-agent system.\\
\textit{Concrete Implementations:} PromptAgent, WriterAgent, ReviewerAgent, EditorAgent, ImageAgent, DirectorAgent, EvaluationAgent

\textbf{PromptAgent}\\
Validates and enhances user input prompts.\\
\textit{Methods:} process\_prompt(prompt) → (enhanced\_prompt, prompt\_type)

\textbf{WriterAgent}\\
Generates structured story content using Ollama LLM.\\
\textit{Methods:} generate\_story(enhanced\_prompt) → (story\_dict, status)

\textbf{ReviewerAgent}\\
Validates story quality and age-appropriateness.\\
\textit{Methods:} review\_story(story\_dict) → (reviewed\_story, status)

\textbf{EditorAgent}\\
Polishes final story and generates PDF export.\\
\textit{Methods:} edit\_story(story\_dict) → (final\_story, status), export\_pdf(story, filename) → pdf\_path

\textbf{ImageAgent}\\
Generates scene illustrations using Stable Diffusion.\\
\textit{Methods:} generate\_image(prompt\_text, filename) → image\_path

\textbf{DirectorAgent}\\
Orchestrates the entire multi-agent pipeline.\\
\textit{Methods:} create\_story(user\_prompt, generate\_images) → response\_dict

\textbf{StoryAgent}\\
Facade providing simplified interface to DirectorAgent.\\
\textit{Methods:} generate\_story(prompt, generate\_images) → (result, status)

\clearpage
